% Generated draft: requires mathematical and source review.
\documentclass[11pt]{article}
\usepackage[margin=25mm]{geometry}
\usepackage{fontspec}
\setmainfont{Latin Modern Roman}
\usepackage{amsmath,amssymb,mathtools,unicode-math}
\setmathfont{Latin Modern Math}
\usepackage{enumitem,needspace,xurl,hyperref}
\setlength{\parindent}{0pt}
\setlength{\parskip}{5pt}
\setlength{\emergencystretch}{3em}
\newcommand{\N}{\mathbb{N}}
\newcommand{\Z}{\mathbb{Z}}
\newcommand{\Q}{\mathbb{Q}}
\newcommand{\R}{\mathbb{R}}
\newcommand{\C}{\mathbb{C}}
\newcommand{\F}{\mathbb{F}}
\newcommand{\A}{\mathbb{A}}
\newcommand{\PP}{\mathbb{P}}
\newcommand{\E}{\mathbb{E}}
\newcommand{\eps}{\varepsilon}
\newcommand{\dd}{\,\mathrm{d}}
\begin{document}
\section*{UnsolvedMath ID: 30000770}
\textbf{Source problem code:} catalogue entry
\subsection{Definitions and mathematical statement}
Let $X$ be a smooth projective algebraic $K(\pi,1)$ space in characteristic zero, as in the source, and put $\omega_X=\det\Omega_X^1$. In the complex topological interpretation, $K(\pi,1)$ means that $X^{\mathrm{an}}$ has fundamental group $\pi$ and $\pi_i(X^{\mathrm{an}})=0$ for $i\geq2$; the extract does not specify the ground field or an alternative algebraic interpretation. For a nef line bundle $L$, its nef reduction is the usual almost-holomorphic reduction with $L$ numerically trivial on general fibers and positive on curves through a very general point that are not contracted. The Iitaka fibration of $\omega_X$, when defined, is the fibration represented by sufficiently divisible pluricanonical systems; the $m$th pluricanonical map is the rational map defined by $H^0(X,\omega_X^{\otimes m})$. Abundance is used in the source as an assumption for the pluricanonical construction; in its usual smooth minimal setting it asserts semiampleness of the nef canonical bundle.\par The source-verification note explicitly says that existence of the indicated finite \'etale cover, making the almost-abelian fibration an abelian scheme, is already a theorem in the report. This existence is therefore retained as reported input, not posed again as the conjecture. Write $f':X'\to Y'$ for the resulting abelian scheme after that indicated cover; an abelian scheme is a proper smooth group scheme with geometrically connected abelian-variety fibers. The exact covering diagram and meaning of the preceding almost-abelian fibration are not included in the extract. The conjectural compatibility asks whether, after this cover, the nef reduction of the canonical bundle, the Iitaka fibration (assuming the source's standard Kodaira-dimension-zero conjecture), and suitable pluricanonical maps (assuming abundance) agree and have common base $Y'$ of general type. For a smooth projective base, general type means $\kappa(Y')=\dim Y'$. The original statement adds the stronger assertion that the canonical bundle $\omega_{Y'}$ should be ample; this must be retained even though the clean/current question omits it. General type alone is not being identified with ampleness. The precise meaning of agreement, the allowed pluricanonical indices, the Kodaira-dimension-zero assumption, and the model and regularity of $Y'$ remain unspecified. The supplied 2026-08-21 triage calls this a conjectural package dependent on unresolved abundance-type questions and labels it open.
\subsection{Short English statement}
After the finite \'etale cover whose existence the source reports as known, ask whether the canonical nef reduction, Iitaka fibration, and suitable pluricanonical maps coincide, with base of general type. The original also conjectures that the base's canonical bundle is ample; the stated conditional assumptions must be kept.
\subsection{Sources}
\begin{itemize}\item Local UnsolvedMath registry, permanent ID 30000770, source problem catalogue entry, \emph{Compatibility of Nef Reduction and Iitaka Fibrations}, proposed year 2007; no individual proposer supplied.\item Supplied citation: \emph{Arithmetic and Differential Galois Groups} (2008), Oberwolfach Reports, DOI 10.4171/owr/2007/26, pp.\ 1443--1520: \url{https://doi.org/10.4171/owr/2007/26}. Supplied literature label: Original Oberwolfach report.\end{itemize}
\end{document}
